\documentclass[11pt]{article}
\usepackage[margin=1in]{geometry}
\usepackage{amsmath,amssymb,amsthm}
\usepackage{hyperref}
\newtheorem{theorem}{Theorem}
\newtheorem{lemma}[theorem]{Lemma}
\newtheorem{proposition}[theorem]{Proposition}
\newtheorem{corollary}[theorem]{Corollary}
\theoremstyle{definition}
\newtheorem{definition}[theorem]{Definition}
\theoremstyle{remark}
\newtheorem{remark}[theorem]{Remark}
\newcommand{\Z}{\mathbb{Z}}
\newcommand{\Q}{\mathbb{Q}}
\newcommand{\N}{\mathbb{N}}

\title{Disproof of Conjecture 00000009614:\\
equal Perron eigenvalues do not force isomorphic dimension groups}
\author{jilint777}
\date{2026-10-04}

\begin{document}
\sloppy
\maketitle

\begin{abstract}
Conjecture 00000009614 asserts that the dimension groups of two mixing shifts of
finite type are order-isomorphic \emph{if and only if} their Perron eigenvalues
$\lambda_A,\lambda_B$ generate the same real number field with agreeing unit group
cosets. We refute the ``if'' direction. The full $4$-shift, $A=[4]$, and the edge
shift of $B=\left(\begin{smallmatrix}3&1\\1&3\end{smallmatrix}\right)$ are mixing
and have the same Perron eigenvalue $\lambda_A=\lambda_B=4$, so every reading of
the right-hand side holds. But $\Delta_A\cong\Z[1/4]$ has rank $1$, while
$\Delta_B$ has rank $2$. So the two dimension groups are not even isomorphic as
abstract groups. The disproof is formalized in Lean~4 (core library only).
\end{abstract}

\section{The conjecture and the clause we refute}
\begin{quote}
\emph{Definition:} The dimension group of an SFT is the ordered
$K_0(\Z^nA,\Z^n_+)$, whose ordered isomorphism type is a flow invariant.
\emph{Conjecture:} The dimension groups of two mixing SFTs are order-isomorphic if
and only if the Perron eigenvalues $\lambda_A$ and $\lambda_B$ generate the same
real algebraic number field with agreeing unit group cosets; explicit examples of
equal entropy over different fields are given by an explicit criterion for coprime
Perron numbers.
\end{quote}
The Chinese version says the same. The conjecture is a biconditional. We refute the
implication
\begin{equation}\label{eq:if}
\bigl[\text{$\lambda_A$, $\lambda_B$ generate the same field with agreeing unit cosets}\bigr]
\;\Longrightarrow\;(\Delta_A,\Delta_A^+)\cong(\Delta_B,\Delta_B^+),
\end{equation}
which makes the biconditional, and hence the conjecture, false. We do not use the
side remark about flow invariance, and we do not need the final clause about
``explicit examples''.

\section{Definitions}
Let $M$ be an $n\times n$ nonnegative integer matrix and $X_M$ its edge shift. We
use row vectors, as in Lind--Marcus~\cite{lm}.

\begin{definition}[Dimension group]\label{def:dg}
$\Delta_M=\varinjlim(\Z^n\xrightarrow{M}\Z^n\xrightarrow{M}\cdots)$. Concretely, its
elements are classes $[v,k]$ of pairs $(v,k)\in\Z^n\times\N$, where
\[
(v,k)\sim(w,l)\iff vM^{l+t}=wM^{k+t}\ \text{for some }t\ge0
\]
(think of $[v,k]$ as ``$vM^{-k}$''). Addition is
$[u,k]+[v,l]=[uM^l+vM^k,\,k+l]$, the zero is $[0,0]$, and $-[v,k]=[-v,k]$. The
positive cone is $\Delta_M^+=\{[v,k]: vM^t\in\Z^n_{\ge0}\text{ for some }t\}$.
The shift map is $\delta_M[v,k]=[vM,k]$.
\end{definition}

This is Krieger's dimension group, i.e.\ the ordered $K_0$ in the conjecture. In
\cite[\S7.5]{lm} the same group is realized inside $\Q^n$: if $R_M$ is the eventual
range of $M$ (the image of $\Q^n$ under $M^n$), then
$\Delta_M\cong\{x\in R_M: xM^k\in\Z^n\text{ for some }k\}$ via
$[v,k]\mapsto vM^n(M|_{R_M})^{-n-k}$. We use the direct limit, and in
Section~\ref{sec:er} we repeat the argument in the eventual-range model, so the
result does not depend on which model is used. Both of our matrices are symmetric,
so row and column conventions (that is, $\Delta_M$ versus $\Delta_{M^{\mathsf T}}$)
agree.

\begin{lemma}\label{lem:group}
$\sim$ is an equivalence relation. Addition and negation are well defined, and
$\Delta_M$ is an abelian group.
\end{lemma}
\begin{proof}
Write $\pi_t(v)=vM^t$. Then $\pi_{s+t}=\pi_t\circ\pi_s$ and each $\pi_t$ is
additive. Reflexivity and symmetry are clear. For transitivity, suppose
$uM^{l+t_1}=vM^{k+t_1}$ and $vM^{m+t_2}=wM^{l+t_2}$. Then with $t=l+t_1+t_2$,
\[
uM^{m+t}=(uM^{l+t_1})M^{m+t_2}=vM^{k+t_1+m+t_2}=(vM^{m+t_2})M^{k+t_1}=wM^{k+t}.
\]
If $(u,k)\sim(u',k')$ via $t$, then $(uM^l+vM^k,k+l)\sim(u'M^l+vM^{k'},k'+l)$
via the same $t$. Both sides become $uM^{k'+t+2l}+vM^{k+k'+l+t}$ after
multiplying by $M^{k'+l+t}$, respectively $M^{k+l+t}$. The sum is symmetric up to
$\sim$, so it is also well defined in the second argument. Negation is similar. Two
identities make the group axioms immediate:
$[v,k]=[vM^s,k+s]$, and $[u,k]+[v,k]=[u+v,k]$. With these, any finitely many
elements can be brought to a common level $k$, where addition is addition in
$\Z^n$.
\end{proof}

\begin{definition}
An \emph{order isomorphism} $(\Delta_M,\Delta_M^+)\cong(\Delta_N,\Delta_N^+)$ is a
bijective additive map $f$ with $f(\Delta_M^+)=\Delta_N^+$. A \emph{group
isomorphism} is a bijective additive map. An isomorphism of dimension triples also
requires $f\delta_M=\delta_Nf$.
\end{definition}
Triple isomorphism implies order isomorphism, which implies group isomorphism. We
show that \emph{no group isomorphism} exists, which rules out all three.

A matrix $M\ge0$ is \emph{primitive} if some power $M^k$ ($k\ge1$) is entrywise
positive. For an essential matrix, $X_M$ is mixing if and only if $M$ is primitive
\cite[\S4.5]{lm}. The \emph{Perron eigenvalue} $\lambda_M$ is the spectral radius.
By the Perron--Frobenius theorem it is the unique eigenvalue with a positive
eigenvector, and $h(X_M)=\log\lambda_M$ \cite[\S4.2--4.4]{lm}.

\section{The example}
Let
\[
A=\begin{pmatrix}4\end{pmatrix},\qquad B=\begin{pmatrix}3&1\\1&3\end{pmatrix}.
\]
$X_A$ is the full $4$-shift.

\begin{proposition}\label{prop:same}
$A$ and $B$ are primitive, so $X_A$ and $X_B$ are mixing SFTs. Also
$\lambda_A=\lambda_B=4$, and $h(X_A)=h(X_B)=\log4$.
\end{proposition}
\begin{proof}
Both matrices are entrywise positive. $(1)A=4(1)$ and $(1,1)B=4(1,1)$ with positive
eigenvectors. Also $(1,-1)B=2(1,-1)$, so $\operatorname{spec}B=\{4,2\}$ and
$\chi_B(x)=x^2-6x+8=(x-4)(x-2)$, with spectral radius $4$. The number of paths of
length $k$ is $\sum_{i,j}(A^k)_{ij}=4^k$ and $\sum_{i,j}(B^k)_{ij}=2\cdot4^k$,
because $B$ has all column sums equal to $4$.
\end{proof}

\paragraph{Every reading of the right-hand side holds.} The right-hand side of
\eqref{eq:if} is a condition on the pair $(\lambda_A,\lambda_B)$. Here
$\lambda_A=\lambda_B=4$, so the condition holds under every interpretation that is
satisfied by two equal numbers. In detail:
\begin{itemize}
\item ``Generate the same field'': $\Q(\lambda_A)=\Q(4)=\Q=\Q(\lambda_B)$.
\item ``Agreeing unit group cosets'': with $K=\Q$ and unit group
$U_K=\mathcal O_K^\times=\{\pm1\}$, every reasonable reading holds trivially for equal
numbers. This includes $\lambda_AU_K=\lambda_BU_K$, $\lambda_A/\lambda_B\in U_K$, the
classes of $\lambda_A,\lambda_B$ in $K^\times/U_K$ agreeing, and the same with $U_K$
replaced by any subgroup (units of an order, totally positive units, \dots).
\item A stronger reading that also compares the Perron eigenvectors (the ideal class
of the module spanned by the eigenvector entries, which is relevant for shift
equivalence when the characteristic polynomial is irreducible) is still satisfied. The normalized
eigenvectors $(1)$ and $(1,1)$ both have entries generating the unit ideal of $\Z$,
so both classes are trivial.
\item Equal entropy: $h(X_A)=h(X_B)=\log4$.
\end{itemize}
In the Lean formalization this is a hypothesis: the right-hand side is an
\emph{arbitrary} relation $\mathrm{Cond}$ with $\mathrm{Cond}(x,x)$ for all $x$.

\section{The dimension groups have different ranks}
Write $a\cdot x$ for the $a$-fold sum ($a\in\Z$). Two elements $x,y$ of an abelian
group are \emph{independent} if $ax+by=0$ with $a,b\in\Z$ forces $a=b=0$.

\begin{lemma}[$\Delta_A$ has rank one]\label{lem:rank1}
Any two elements of $\Delta_A$ are dependent. More generally this holds for any
$1\times1$ matrix $M$.
\end{lemma}
\begin{proof}
Let $p=[u,k]$, $q=[w,l]$ with $u,w\in\Z$. Lift both to level $K=k+l$:
$p=[\mu,K]$, $q=[\nu,K]$ with $\mu=uM^l$, $\nu=wM^k\in\Z$. At a common level,
integer combinations are computed on numerators:
$a p+bq=[a\mu+b\nu,K]$. If $\mu=0$ then $1\cdot p+0\cdot q=0$. Otherwise
$(a,b)=(\nu,-\mu)\ne(0,0)$ gives $\nu p-\mu q=[\nu\mu-\mu\nu,K]=0$.
\end{proof}
(In fact $[u,k]\mapsto u/4^k$ is an isomorphism $\Delta_A\cong\Z[1/4]$, with
$\Delta_A^+=\Z[1/4]_{\ge0}$. We do not need this.)

\begin{lemma}[$\Delta_B$ has rank at least two]\label{lem:rank2}
$x=[e_1,0]$ and $y=[e_2,0]$ are independent in $\Delta_B$.
\end{lemma}
\begin{proof}
$ax+by=[(a,b),0]$. If this is $0=[(0,0),0]$, then $(a,b)B^t=(0,0)$ for some $t$. Now
$vB=wB$ means $3v_1+v_2=3w_1+w_2$ and $v_1+3v_2=w_1+3w_2$, hence $v=w$, since
$\det B=8\ne0$. By induction $B^t$ is injective on $\Z^2$, so $(a,b)=(0,0)$.
\end{proof}

\begin{theorem}\label{thm:main}
There is no injective additive map $\Delta_B\to\Delta_A$. In particular $\Delta_A$
and $\Delta_B$ are not isomorphic as groups. A fortiori
$(\Delta_A,\Delta_A^+)\not\cong(\Delta_B,\Delta_B^+)$ as ordered groups, and the
dimension triples are not isomorphic.
\end{theorem}
\begin{proof}
Let $f$ be additive and injective, and let $x,y$ be as in Lemma~\ref{lem:rank2}. By
Lemma~\ref{lem:rank1} there are $(a,b)\neq(0,0)$ with $af(x)+bf(y)=0$, so
$f(ax+by)=0=f(0)$, and hence $ax+by=0$, contradicting Lemma~\ref{lem:rank2}. An
isomorphism $\Delta_A\to\Delta_B$ would have an additive injective inverse.
\end{proof}

\begin{corollary}
Implication \eqref{eq:if} is false for $A=[4]$ and
$B=\left(\begin{smallmatrix}3&1\\1&3\end{smallmatrix}\right)$, under every reading of
its hypothesis that holds when $\lambda_A=\lambda_B$, and for every notion of
isomorphism (group, ordered group, dimension triple). Therefore Conjecture
00000009614 is false.
\end{corollary}

\section{Robustness}\label{sec:er}
\paragraph{Eventual-range model.} In the model of \cite[\S7.5]{lm} the argument is
even more direct. $R_A=\Q$ and $\Delta_A=\{x\in\Q:4^kx\in\Z\text{ for some
}k\}=\Z[1/4]$. Any two elements $m/4^k$ and $m'/4^{k'}$ satisfy
$m'4^k\cdot\frac{m}{4^k}-m4^{k'}\cdot\frac{m'}{4^{k'}}=0$. Since $\det B=8\ne0$, we
have $R_B=\Q^2$ and $\Z^2\subseteq\Delta_B\subseteq\Q^2$. So
$\operatorname{rank}\Delta_B=\dim_\Q(\Delta_B\otimes\Q)=2\neq1=\operatorname{rank}\Delta_A$.
Rank is an invariant of abstract abelian groups.

\paragraph{A second obstruction (dimension triples).} By Krieger's theorem
\cite[\S7.5]{lm}, $(\Delta_A,\Delta_A^+,\delta_A)\cong(\Delta_B,\Delta_B^+,\delta_B)$
if and only if $A$ and $B$ are shift equivalent. Shift-equivalent matrices have the
same nonzero spectrum \cite[\S7.4]{lm}. Here the nonzero spectra are $\{4\}$ and
$\{4,2\}$. Equivalently, $\operatorname{tr}A^k=4^k\neq4^k+2^k=\operatorname{tr}B^k$,
which counts the period-$k$ points. This independent argument covers the triple
reading. The rank argument covers all readings.

\paragraph{Mixing.} Both matrices are positive, so the shifts are mixing under any
presentation (edge shift; $X_A$ is also the full $4$-shift as a vertex shift).

\paragraph{Why the conjecture fails.} The rank of $\Delta_M$ is the number of nonzero
eigenvalues of $M$ (with algebraic multiplicity). Any condition that depends only
on $\lambda_M$ ignores the non-Perron part of the spectrum. Our $B$ has reducible
characteristic polynomial. The conjecture places no restriction on the
characteristic polynomial, and we do not claim anything about a variant restricted
to irreducible characteristic polynomials.

\section{Lean formalization}
The directory \texttt{lean4/} is a self-contained Lean~4.19.0 project (core library
only, no Mathlib). File \texttt{Main.lean}, namespace \texttt{DimGroup}:
\begin{itemize}
\item Vectors are \texttt{Fin n -> Int} and matrices are \texttt{Fin n -> Fin n -> Int}.
\texttt{rmul M v} is $vM$ (row vector) and \texttt{pw M t v} is $vM^t$, for
arbitrary $n$.
\item \texttt{Rel M} is the relation of Definition~\ref{def:dg}. Its equivalence is
proved, and \texttt{DG M} is the quotient type $\Delta_M$. \texttt{Add},
\texttt{Neg} and \texttt{Zero} are defined on the quotient, and well-definedness is
proved. \texttt{add\_comm'}, \texttt{add\_assoc'}, \texttt{zero\_add'} and
\texttt{neg\_add'} prove the abelian group axioms. \texttt{Pos M} is the positive
cone.
\item \texttt{GroupIso} (additive bijection) and \texttt{OrderIso} (additionally
$x\in\Delta^+\iff f(x)\in\Delta^+$) are structures. \texttt{OrderIso.refl} shows
these notions are not vacuous.
\item \texttt{rank\_one}: for every $1\times1$ matrix, any $p,q$ satisfy
$(a+1)p+(b+1)q=(c+1)p+(d+1)q$ for some $(a,b)\ne(c,d)$. This is the integer
relation with its positive and negative parts separated, so that only positive
multiples \texttt{nsm1} are needed. \texttt{rank\_two}: in $\Delta_B$, such an
identity for $[e_1,0],[e_2,0]$ forces $(a,b)=(c,d)$.
\texttt{no\_injective\_hom}, \texttt{not\_groupIso\_AB} and
\texttt{not\_orderIso\_AB} prove Theorem~\ref{thm:main}.
\item \texttt{Primitive M} means $M\ge0$ and some $M^k$, $k\ge1$, is positive.
\texttt{IsPerronEig M lam} means there is a positive $v$ with $vM=\lambda v$. These
are proved for $A,B$ with $\lambda=4$ (\texttt{primitive\_A}, \texttt{perron\_B},
\dots). \texttt{second\_eig\_B} gives the eigenvalue $2$, and
\texttt{pathCount\_A}, \texttt{pathCount\_B} give the path counts $4^k$ and
$2\cdot4^k$.
\item \texttt{IfDirection Cond} is the statement ``for all $n,m$ and all
$A\in\Z^{n\times n}$, $B\in\Z^{m\times m}$ primitive with integer Perron
eigenvalues $\lambda_A,\lambda_B$ satisfying $\mathrm{Cond}(\lambda_A,\lambda_B)$,
$\Delta_A$ and $\Delta_B$ are order-isomorphic''. This is a special case of the
``if'' direction. \texttt{conjecture\_00000009614\_false} proves
$\neg$\texttt{IfDirection Cond} for every reflexive \texttt{Cond}.
\texttt{conjecture\_00000009614\_iff\_false} refutes the biconditional,
\texttt{conjecture\_00000009614\_group\_false} refutes the version with mere group
isomorphism, and \texttt{conjecture\_00000009614\_false\_Q} instantiates the
concrete reading $\lambda_A=\pm\lambda_B$ (field $\Q$, units $\{\pm1\}$).
\texttt{nonvacuous} records that all hypotheses are satisfiable.
\end{itemize}
Not formalized: the Perron--Frobenius theorem itself (that a positive eigenvector
characterizes the spectral radius, used only to identify \texttt{IsPerronEig} with
``Perron eigenvalue''; for $B$ the full spectrum $\{4,2\}$ is exhibited by
eigenvectors in Lean), the identification of the direct limit with the
eventual-range model, and Krieger's theorem (used only in the second, redundant
obstruction). \texttt{lake build} succeeds with no warnings. There is no
\texttt{sorry}, no \texttt{native\_decide}, and no added axiom. \texttt{\#print
axioms} shows only \texttt{propext} and \texttt{Quot.sound}.

\section{Reproduction}
\begin{verbatim}
cd lean4 && lake build     # Lean 4.19.0, core only
python3 verify.py          # Python 3 standard library
pdflatex report.tex
\end{verbatim}
\texttt{verify.py} checks the same facts independently, with exact rational
arithmetic. It computes the characteristic polynomials, the spectral radius (also by
power iteration), and the path counts. It tests rank one and rank two in the
eventual-range model, and runs a brute-force check of the identifications in the
direct limit. It also verifies the trace and spectrum invariants.

\begin{thebibliography}{9}
\bibitem{lm} D.~Lind and B.~Marcus, \emph{An Introduction to Symbolic Dynamics and
Coding}, Cambridge University Press, 1995 (\S4.2--4.5 Perron--Frobenius theory and
mixing; \S7.4 shift equivalence; \S7.5 dimension groups and Krieger's theorem).
\bibitem{kr} W.~Krieger, On dimension functions and topological Markov chains,
\emph{Invent. Math.} 56 (1980), 239--250.
\end{thebibliography}
\end{document}
